\documentclass[../../../uem_tok_q.tex]{subfiles}

\begin{document}
    $n$を2以上の整数とする．1から$n$までの数字が書かれた札が各1枚ずつ合計$n$枚あり，%
    横一列におかれている．1以上$(n-1)$以下の整数$i$に対して，%
    次の操作$(\mathrm{T}_i)$を考える．
    \begin{itemize}[label={$(\mathrm{T}_i)$}, leftmargin=4\zw, labelsep=.7\zw, topsep=4pt]
        \item 左から$i$番目の札の数字が，左から$(i+1)$番目の札の数字よりも大きければ，%
            これら2枚の札の位置を入れかえる．そうでなければ，札の位置をかえない．
    \end{itemize}

    最初の状態において札の数字は左から$A_1,\:A_2,\:\dots,\:A_n$であったとする．%
    この状態から$(n-1)$回の操作%
    $(\mathrm{T}_1),\:(\mathrm{T}_2),\:\dots,\:(\mathrm{T}_{n-1})$を順に行った後，%
    続けて$(n-1)$回の操作$(\mathrm{T}_{n-1}),\:\dots,\:(\mathrm{T}_2),\:(\mathrm{T}_1)$%
    を順に行ったところ，札の数字は左から$1,\:2,\:\dots,\:n$と小さい順に並んだ．%
    以下の問いに答えよ．

    \begin{enumerate}
        \item $A_1$と$A_2$のうち少なくとも一方は2以下であることを示せ．
        \item 最初の状態としてありうる札の数字の並び方$A_1,\:A_2,\:\dots,\:A_n$%
            の総数を$c_n$とする．$n$が4以上の整数であるとき，$c_n$を$c_{n-1}$と$c_{n-2}$を用いて表せ．
    \end{enumerate}
\end{document}
